% 【本文件写什么】api-v0 契约层（叶子，只写语义不写实现）
%   - crate 路径：os/components/wateros-ipc/ipc-signal/signal-api/api-v0/src/
%   - 列出并说明：pub trait、核心类型、错误枚举、常量
%   - 每个公开 API 的前置条件、后置条件、线程/中断上下文限制
%   - 与 Linux/用户态约定的对齐点与 deliberate 差异
%   - v0 版本当前行为与后续可替换 extension 点
% 【事实来源】
%   - 上述 crate 下全部 .rs 源文件
%   - docs/exports/public-api/ 中对应符号表（以源码为准）
% 【不写什么】
%   - impl 内部数据结构、汇编、具体算法步骤

\subsubsection{signal-api-v0}

\paragraph{信号编号与 SA 标志。}
导出\code{NSIG}（64）、常用\code{SIG*}常量、\code{SIG\_DFL}/\code{SIG\_IGN}、\code{sigprocmask}的\code{SIG\_BLOCK}/\code{SIG\_UNBLOCK}/\code{SIG\_SETMASK}，以及\code{SA\_NODEFER}、\code{SA\_RESETHAND}、\code{SA\_RESTART}等\code{rt\_sigaction}标志。\code{ITIMER\_REAL}/\code{ITIMER\_VIRTUAL}/\code{ITIMER\_PROF}对应三类 interval timer。

\paragraph{核心类型。}
\code{SignalSet}为 64 位掩码，提供\code{contains}/\code{insert}/\code{union}/\code{first\_signal}等操作。\code{SignalAction}为\code{repr(C)}，字段\code{handler}、\code{flags}、\code{restorer}、\code{mask}与 Linux\code{struct sigaction}布局对齐。\code{PendingSignal}描述 trap 交付帧：信号号、disposition 快照与进入 handler 前应恢复的\code{previous\_mask}。

\paragraph{投递分类。}
\code{SignalDelivery}枚举\code{Ignored}、\code{Pending}、\code{Terminate}、\code{Stop}、\code{Continue}；\code{SignalDispatch}附带可选\code{target\_task\_id}供 syscall 唤醒阻塞线程。\code{IntervalTimerSpec}以纳秒表示间隔与剩余值。

\paragraph{辅助契约。}
\code{valid\_signal}、\code{immutable\_signal(SIGKILL)}、\code{default\_ignored}/\code{default\_terminates}描述默认 disposition 规则；\code{SignalError}覆盖非法信号、非法\code{how}、未知任务/进程与非法 timer 种类。状态机实现位于\code{ipc-signal}聚合\code{lib.rs}，不在本 api crate。
